% Run GCT_KRON_20261120; GPT-6 Astra Pro; effort not exposed.
\section{An optimized differential recurrence and a general padded-degree gate}

\subsection*{A-022. A uniform degree restriction for actual multiplicity obstructions}
For $L\geq1$ define integers $N_a(L)$, $a\geq2$, by
\[
 N_2(L)=1,\qquad s=\min(a-1,L),\qquad
 N_a(L)=s+\max\{N_{a-1}(L),\min((a-1)s-1,L)\}
                       \quad(a\geq3).                 \tag{A-022.1}
\]
They are nondecreasing in $L$. For every complex vector space
$V$, every $b\geq1$, and every $\lambda\vdash ab$ with
$\lambda_2\leq L$,
\[
 b\geq N_a(L)
 \quad\Longrightarrow\quad
 I(\operatorname{Ch}_b(V))_a[\lambda]=0.                \tag{A-022.2}
\]
The bounds satisfy
\[
 N_a(L)\leq B_a(L),\qquad
 N_a(L)\leq a(a-1)-1,                                 \tag{A-022.3}
\]
where $B_a$ is A-021's bound. This recurrence uses both
possible injectivity tests at the last differential step. It
is an optimization within those degree and second-row tests,
not an assertion of the least possible Chow stabilization bound.

For the GCT application let $n\geq m\geq1$, $d\geq2$.
If
\[
                              n\geq N_d(dm),           \tag{A-022.4}
\]
then $q^m_\lambda(d[n])\leq K_n(\lambda)$ for every partition
$\lambda\vdash nd$. More precisely, every type with
$q^m_\lambda(d[n])>0$ has full ambient determinant multiplicity.
Consequently an actual multiplicity obstruction requires
\[
 n<N_d(dm)\leq
 \min\left\{d(d-1)-1,\;dm+\binom d2-1\right\}.        \tag{A-022.5}
\]
In particular the broad condition
$n\geq dm+\binom d2-1$ already excludes all actual
multiplicity obstructions in that degree, by A-021 alone.

There is a closed finite maximum for the sharper bound needed
in (A-022.4): since $L=dm\geq d$,
\[
 \boxed{\quad N_d(L)=\binom d2+
 \max\left(0,\max_{2\leq j\leq d-1}
      \left\{\min(j^2-1,L)-\binom j2\right\}\right).\quad}
                                                               \tag{A-022.6}
\]
The empty inner maximum gives $N_2(L)=1$. For $d\geq3$
only the two indices
$\min(d-1,\lfloor\sqrt{L+1}\rfloor)$ and
$\min(d-1,\lfloor\sqrt{L+1}\rfloor+1)$ need be inspected.
An explicit sufficient bound, slightly weaker than the exact
maximum, is
\[
 n\geq \binom d2+
 \left\lfloor\frac{dm-1+\lfloor\sqrt{dm+1}\rfloor}{2}
                                                    \right\rfloor.
                                                               \tag{A-022.7}
\]
Thus the exact degree threshold can be computed using only
integer arithmetic and two candidates.

\begin{proof}[Proof of the recurrence]
The two applicable cases of A-020's differential kernel lemma
say that $D_{Y,U}$ is injective on type $\lambda$ whenever
the source degree $p$ satisfies
\[
                            p>\min(q,\lambda_2),       \tag{A-022.8}
\]
where $q$ is the destination degree. Indeed $p>q$ gives
injectivity on the entire space, while $p>\lambda_2$ gives
injectivity on the specified type with all spectators.

We prove (A-022.2) for each fixed upper bound $L$, by induction
on $a$. One-row types have no Chow equations by A-021's
one-row argument. At $a=2$ there are no quadratic equations
for any $b\geq1$, again as in A-021. Now set $r=a-1$,
$s=\min(r,L)$ and $m=b-s$. The recurrence ensures
\[
       m\geq N_r(L)\geq1,
       \qquad m\geq\min(rs-1,L).                      \tag{A-022.9}
\]
For $T$ of type $\lambda$, use the same first sequence as
in A-021, starting from
$H=T(x^sQ,\ldots,x^sQ,t^b)$, $Q=y_1\cdots y_m$.
Its iterated derivative is $b!/s!$ times
$T(x^sQ,\ldots,x^sQ,t^sQ)=0$, from binary factorization.
At every individual step the source degree is at least
$s+1$, and the untouched destination has degree $r$.
If $s=r$, it is injective by $p>q$. If $s=L$, it is
injective by $p>L\geq\lambda_2$. The images retain type
$\lambda$, so $H=0$.

Keeping $x,t$ as polynomial spectators gives the equivariant
image in
\[
 \operatorname{Sym}^{rs}V^*\otimes\operatorname{Sym}^bV^*
                              \otimes I(\operatorname{Ch}_m(V))_r.
\]
The coefficientwise ideal argument is the one proved in A-021.
By $m\geq N_r(L)$ and the induction hypothesis, the last
factor contains no type $\nu$ with $\nu_2\leq L$.
Two horizontal Pieri steps cannot decrease its second row.
Hence the image has no type $\lambda$, and
$T(x^sQ_1,\ldots,x^sQ_r,t^b)=0$ for independent $Q_i$.

Apply the second sequence $D_{x,v_i}^s$, $i=1,\ldots,r$,
to $T(v_1^b,\ldots,v_r^b,t^b)$. The resulting polynomial
is zero by the preceding identity. Before every individual
derivative, $p\geq m+1$ and $q\leq rs-1$.
By (A-022.9) and $\lambda_2\leq L$,
\[
 p>\min(rs-1,L)\geq\min(q,\lambda_2).
\]
Thus (A-022.8) makes every step injective on the relevant
type, and $T$ vanishes on all tuples of pure powers. Spanning
and multilinearity give $T=0$, proving (A-022.2).

For monotonicity, $s=\min(a-1,L)$ and
$\min((a-1)s-1,L)$ are nondecreasing in $L$; induction
in (A-022.1) proves the assertion. Also
$\min((a-1)s-1,L)\leq L\leq B_{a-1}(L)$ and
$N_{a-1}(L)\leq B_{a-1}(L)$ by induction. Adding $s$
therefore gives $N_a(L)\leq B_a(L)$. Finally $s\leq r$
and $\min(rs-1,L)\leq r^2-1$, while
$N_r(L)\leq r(r-1)-1\leq r^2-1$ by induction.
Thus $N_a(L)\leq r+r^2-1=a(a-1)-1$, proving (A-022.3).
\end{proof}

\begin{proof}[Proof of the padded-degree consequence]
We include the needed padding weight restriction. The orbit
closure of $\ell^{n-m}\operatorname{per}_m$ is contained in the
closure of all forms $\ell^{n-m}f$, with $f$ any degree-$m$
form. Pulling back a degree-$d$ polynomial along this
multiplication morphism embeds the degree-$d$ coordinate
module of that larger family into
\[
          \operatorname{Sym}^{d(n-m)}V\otimes
                         \operatorname{Sym}^d(\operatorname{Sym}^mV).
\]
The first factor records the coefficients of $\ell$, and the
second those of $f$. If a type $\lambda$ occurs on the
permanent orbit closure, it must occur in this module.
Decompose the second factor into $S_\nu V$ with $|\nu|=dm$.
Horizontal Pieri says that $\lambda/\nu$ is a horizontal
$d(n-m)$ strip. In particular $\lambda_2\leq\nu_1\leq dm$.
Also, a horizontal strip has at most one added box per column,
so $\lambda_1\geq d(n-m)$ and $|\bar\lambda|\leq dm$.
This is a statement about every
positive permanent-side multiplicity, not only a particular
highest-weight vector.

Now assume $n\geq N_d(dm)$. Apply (A-022.2) with $L=dm$
to every type with $q^m_\lambda>0$. Its Chow-ideal component
is zero. Since products of $n$ linear forms are diagonal
determinants, their cone lies in the determinant orbit closure;
therefore every such determinant-ideal component is zero as
well. This gives $K_n(\lambda)=a_\lambda(d[n])\geq q^m_\lambda$.
For $q^m_\lambda=0$ the comparison is automatic. Thus no
actual multiplicity obstruction occurs in this degree.
For $L=dm\geq d$, (A-021.2) reads
$B_d(dm)=dm+\binom d2-1$, and (A-022.3) gives (A-022.5).

The simpler degree restriction can be written without $B,N$:
every actual obstruction has
\[
 d>
 \frac{-(2m-1)+\sqrt{(2m-1)^2+8(n+1)}}{2}.              \tag{A-022.10}
\]
Indeed $n<dm+d(d-1)/2-1$ is exactly the strict quadratic
inequality whose positive root is displayed. This is a
necessary condition; the exact recurrence is stronger.
\end{proof}

\begin{proof}[Proof of the finite maximum]
When $L\geq d$, every $s$ occurring for $a\leq d$ is $a-1$.
Unroll (A-022.1) from $N_2=1$. The sum of the additive
increments $2+\cdots+(d-1)$ is $\binom d2-1$.
An intermediate term $\min(j^2-1,L)$ is inserted just
before adding $j$, so the increments before its insertion
sum to $\binom j2-1$. Taking the maximum of the initial
term and all inserted terms yields precisely (A-022.6).

Put $k=\lfloor\sqrt{L+1}\rfloor$. For integer $j\leq k$,
the summand in braces is $\binom{j+1}{2}-1$, increasing
with $j$. For $j\geq k+1$ it is $L-\binom j2$, decreasing
with $j$. Hence the maximum on $2\leq j\leq d-1$ occurs
at the two stated, clipped indices. The case $d=2$ was
already specified.

For completeness, the corresponding real function of $j$
has its maximum at $j=\sqrt{L+1}$, with value
$(L-1+\sqrt{L+1})/2$. Its integer maximum is thus at most
$\lfloor(L-1+\sqrt{L+1})/2\rfloor$, which equals
$\lfloor(L-1+\lfloor\sqrt{L+1}\rfloor)/2\rfloor$.
The equality follows by separating the integer and fractional
parts of the square root: its fractional part divided by two
cannot cross the next integer from an integer or half-integer.
Adding $\binom d2$ proves (A-022.7).
\end{proof}

\paragraph{Scope and geometric dictionary.}
This is a restriction on actual orbit-closure evaluation ranks.
For $m$ small compared with $d$, the sufficient expression
(A-022.7) has leading terms $\binom d2+dm/2$, with its
displayed square-root correction. The exact expressions (A-022.6)--(A-022.7)
are used for conclusions. No conclusion about all higher
degrees, global Kronecker saturation, or $\mathsf P$ versus
$\mathsf{NP}$ follows from this reduction alone.

The Chow and padding multiplication maps are morphisms of
affine schemes over $\mathbb Z$, and their homogeneous
coordinate modules are comodules as in SGA 3, I,
Propositions 4.7.2--4.7.3, pp.~34--36. Pieri is Macdonald,
I, (5.16), p.~73. The symmetric and Schur functors are
interpreted in the exact $K_0$ framework of SGA 6, VI,
Theorem 3.3, pp.~392--393. The proof uses the reduced
characteristic-zero fibers and preserves all isotypic
multiplicity copies; integral definitions alone do not
assert the same kernel statement in every characteristic.
